% !TeX root = ../papers/nc-p3.tex
\section{整な非正規八次曲面からの反対根の復元}
\label{sec:cofactor-curve-reconstruction}
有限窓の一配置 $X$ について $W=H^0(\mathcal I_X^2(6))$、その評価による有理写像を
$\phi:\mathbb P^3\dashrightarrow\mathbb P(W^*)$ とする。以下で根は実際の反対根、
または一般階数4で全階数平方を与えるものを数える。正規化の行符号は除く。
\begin{theorem}\label{thm:cofactor-curve-reconstruction}
$X$ を三重に含む八次式がスカラーを除いて一つで、その曲面 $D=V(d)$ が整とする。
次表の内在条件の下で、恒等根以外の全ての反対根を有限個の線束と行列から復元できる。
$\Gamma=\overline{\phi(D)}$ とし、表の曲線は表示された線型系で埋め込まれているとする。
\[
\begin{array}{c|l|c}
\dim W&\Gamma&\text{非恒等根数の上界}\\\hline
7&\mathbb P^1,\ |\mathcal O(6)|&1\\
8&\text{滑らかな種数1曲線の完全八次系}&4\\
8&\text{節点を持つ整な算術種数1曲線の完全八次系}&3\\
8&\text{尖点を持つ整な算術種数1曲線の完全八次系}&2
\end{array}
\]
とくに整な反対八次行列の余因子が $W$ を張り、その次元が7か8なら、
この表のいずれかが適用される。七余因子正規形の一般の三割平面対は最初の行を満たす。
\end{theorem}
\begin{proof}[核点の線束]
$N$ が非恒等な実際の反対根であれば行列式は非零であり
（定理\ref{thm:all-finite-opposite-octics}）、定数倍を除き $d$ である。
十個の六次余因子は $W$ に入り、$\Gamma$ から $\mathbb P^9$ への線型写像を与える。
$D$ の一般点の階数は3なのでこの像はVeronese上にあり、一般核点は三空間を線型に張る。
実際それらを消す線型形式があれば、対応する六次余因子の四つの線型結合が
$D$ で消えて恒等零になり、一般可逆性に反する。
$\Gamma$ の正規化上で線型写像の固定因子 $B$ を除き、Veroneseの逆で得る
三空間への射の線束を $L$ とすれば
\[
             L^{\otimes2}(B)=\nu^*\mathcal O_\Gamma(1).
\]
$L$ には基点なく非退化な四切断がある。六次の有理正規曲線では
$\deg L\ge3$ なので $B=0$、四切断は $H^0(\mathcal O(3))$ の全体である。
従って出力の基底変更を除いて表示はただ一つとなる。
滑らかな種数1の八次曲線でもRiemann--Rochにより $\deg L\ge4$ で、
$B=0$、四切断は完全系となる。$L^2=\mathcal O_\Gamma(1)$ の解は二捩れの四通りである。

特異な算術種数1の場合、正規化は射影直線で $\deg L=3,4$ である。
4なら $B=0$ だから線型写像の中心は $\Gamma$ を避け、そのVeronese逆も
特異曲線上の射となる。線束は $\Gamma$ 上で平方根となり、四切断は全てである。
実際双対化層は自明で、正次数の線束に第一コホモロジーはない。
一般化Jacobi多様体は節点で乗法群、尖点で加法群なので二捩れはそれぞれ二個、一個である。
これは正規化で節点両枝の繊維を貼る非零定数、または尖点の一次近傍の単元を
定数単元で割った一次係数を使っても分かる。
3なら正規化上の三次完全系で、その平方の線型包は七次元である。
従って $\mathbb P^7$ からの線型射影の中心は一点である。
取り除く固定因子の次数は2なので、その中心は曲線上の特異点でなければならない。
滑らかな点ではその点を通る全超平面系の引戻し固定因子は単純一点である。
特異点は一つしかなく、中心を指定すれば次数6の像とその三次表示は基底を除いて一意である。

実際の行列で余因子が $W$ を張れば、$\Gamma$ は定理\ref{thm:seven-cofactor-trisecant-grassmannian}
の核曲線の二次Veroneseである。三次の場合は七切断全て、整な二次完全交叉四次曲線の場合は
二次式の制限の八切断全てとなる。後者が特異なら平面切断体内の曲線の特異性を仮定する必要はない。
種数差が1であるため特異点一つで $\delta=1$ である。
正規化との余長1の局所環は二つの滑らかな枝の定数が一致する環か、
一枝で一次係数が零な環であり、それぞれ節点、尖点である。
\end{proof}

\begin{proof}[有限候補の行列と標本]
基底の積から $\Gamma$ 上でVeroneseを与える十個の一次座標を作る。
完全八次系では四次線束の切断の積を $H^0(\mathcal O_\Gamma(1))$ と比較し、
特異点から射影する場合は中心を通る七切断を使う。
それらの一次式に $W$ の六次式を代入した対称行列を $P(x)$ とする。
$D$ の一般点で階数1なので、素元 $d$ の離散付値環で単元一成分を消去すると
\[
       d^2\mid\operatorname{adj}P,\qquad \det P=c d^3
\]
となる。最後は次数24が等しいことによる。$c=0$ は除き、
$N=\operatorname{adj}P/d^2$ は二次式行列で
$\det N=c^3d$、$\operatorname{adj}N=c^2 P$ となる。
十成分が二次式の基底か、$X$ 上で階数1で基点がないか、得た座標の三次評価は同型かを検査する。
階数1なら各局所Artin因子で単元対角の平方根を取って行列を列の平方に分解できる。
これで条件を通る平方根を全て得る。もとの行列の六次余因子表示は
$\Gamma$ が全空間を線型に張ることから、この構成で復元される。
関係六空間を持つ辺になることと正則性は、その根を使う窓でさらに必要となる。

一般の三割平面対では、$T$ が定める二四次式の共通根を印すと
$\mathbb P^1$ の各点上の二本の独立一次式で $x$ が切られ、射影直線束になる。
独立性は二本の四次曲線が交わらないことによる。この束の像が八次曲面である。
一般の $T$ では一般最大公約次数はそこで1で、行列式は重複しない。
実際 $\operatorname{Gr}(2,S_4)$ で二次共通因子を持つ鉛筆は
$\mathbb P(S_2)\times\operatorname{Gr}(2,S_2)$ の像で次元4以下である。
一次共通因子が一つで単純な開集合は滑らかな結果式因子である
（共通根の評価式を根パラメータで消去する）。
推移群による一般移動で三次元の鉛筆族との交叉は前者で次元1以下、
後者で横断的である。これは群作用写像の滑らかさと標数零の一般滑らかさによる。
従って $D$ は整で七余因子が存在する。
\path{computations/verify_cofactor_curve_reconstruction.py} は既出の整数 $T$ の
$x$ 側で全一階偏導と二階偏導の評価核を計算し、
$\dim H^0(\mathcal I_X^2(6))=7$、$\dim H^0(\mathcal I_X^3(8))=1$
の上界を良い素数で検証する。一般点での逆の不等式は実在する七余因子と行列式による。
\end{proof}

\begin{proposition}\label{prop:elliptic-cofactor-quartic-pencil}
滑らかな種数1の行で一つの全階数反対根が実在するとする。
四つの二捩れに対応する $\mathbb P(W^*)$ 上の対称行列式は、一本の四次鉛筆の
相異なる四つの元である。そのうち少なくとも三つは、前定理で恒等零にならず二次式行列を与える。
$X$ が被約で $D$ の各 $X$ 点における重複度がちょうど3なら、これらの二次式行列は
そこで全て階数1である。十成分の一次独立性や正則窓への延長は別に検査する。
\end{proposition}
\begin{proof}
$E=\Gamma$ の原点 $o$ を $\mathcal O_E(1)=\mathcal O_E(8o)$ と選び、
$L_a=\mathcal O_E(3o+a)$、$a$ の群逆元を $-a$ とする。
$L_aL_{-a}=\mathcal O_E(8o)$ の切断の積に $b\in W^*$ を作用させて
四行の対合を得る。二捩れで対称となり、$a,-a$ は互いに転置を与える。
底 $E$ 上でこの行列式の係数が動く線束は次数2である。
実際 $E\times E$ の対角 $\Delta$ により
$\mathcal L=\mathcal O(\Delta)\otimes p_1^*\mathcal O(3o)\otimes p_2^*\mathcal O(-o)$
と正規化すると、$p_{2*}\mathcal L$ は階数4、次数 $-1$ である。
対角への制限で $\mathcal O(-o)^3$ と $\mathcal O(2o)$ の拡張となるからである。
底で $a\mapsto-a$ とした束との積は $p_1^*\mathcal O(8o)$ である。
各ファイバーでの同型を割った線束は底からの引戻しであり、
$\{o\}\times E$ 上の正規化で自明と分かる（自明なファイバーの大域切断と底変換でもよい）。
行列式は二つの階数4束の双対行列式の積に値を取り、この次数2の線束の切断は二次元である。

どの $a$ でも行列式は非零である。四点での評価の一般一次結合で、
二つの四切断系の両方の評価が可逆となるように点を選べるためである。
さらに異なる四点の和が $A$ であるとき、その評価の全係数が非零な結合での行列式は、
$\mathcal O(A)=L_a$ または $L_{-a}$ のときちょうど零である。
積の評価が二つの四次評価行列と可逆対角行列との積だからである。
各四次線形系には相異なる四点の因子が存在するので、$a$ を符号まで動かすと
相異なる四次式を得る。これで鉛筆と四つの元の相違が従う。
$\phi$ での代入はこの二次元空間を引き戻し、その四つの元は $d^3$ の定数倍である。
既知の反対根により像は零でなく、鉛筆で恒等零になる元は一つだけである。

$P$ の成分は $I_X^2$ に入るので、$c\ne0$ の行列 $N$ の全三次小行列式も
各点 $p\in X$ で二重に消える。階数2なら単元二行を消去した残り二行の成分は
全て二重に消え、$d$ の位数は4以上となる。階数0でも同じく行列式は四重に消える。
階数3以上は余因子の消滅に反する。これで最後の主張を得る。
付属計算では $y^2=x(x-1)(x-3)$ で四つの積写像を明示し、四次式の各対が
実際に独立であることを有理数上でも確認する。
\end{proof}
